
\chapter{Tensorgrad}

Throughout this book we have manipulated tensor diagrams by hand.
The \emph{tensorgrad} library (\url{https://github.com/thomasahle/tensorgrad})
does the same manipulations symbolically in Python: tensors are represented as
graphs with \emph{named} edges, so there is no index bookkeeping, no
transpose-juggling, and no einsum strings to keep consistent.
On top of this representation the library implements symbolic differentiation,
diagram simplification, and compilation to numeric PyTorch code.
In this chapter we give a short overview of the design, and of the problems a
tensor-diagram engine has to solve that a classical computer-algebra system does not.

\section{Representation}

Every expression is a subclass of \texttt{Tensor}.
The leaves are \texttt{Variable} (a named tensor whose dimensions are sympy
symbols) and the constants \texttt{Delta} (the copy/identity tensor, which also
doubles as a hyper-edge of any order) and \texttt{Zero}.
Composite expressions are built from \texttt{Sum}, \texttt{Product},
\texttt{Function}, and \texttt{Derivative}; a lightweight \texttt{Rename}
node relabels edges lazily.
Nonlinear functions such as \texttt{exp}, \texttt{log}, and \texttt{softmax}
are described by a \texttt{FunctionSignature}, which declares which input edges
are consumed, which output edges are produced, and how to differentiate---so
users can add new functions without touching the core.

The key decision is that edges are identified by \emph{name}, not by position.
A matrix is a tensor with two edges, say $i$ and $j$, and contraction happens
simply by forming a \texttt{Product} of tensors that share an edge name.
There is no distinction between a matrix and its transpose beyond which names
its edges carry, which matches how tensor diagrams themselves work.
The price is that renaming becomes a first-class operation the engine must
handle carefully---for instance when a derivative introduces new edges that
would otherwise clash with existing names.

\section{Canonicalization}

The central algorithmic problem is deciding when two diagrams are \emph{the same}.
Since a diagram is a labeled multigraph, this is exactly graph isomorphism:
each expression can produce a NetworkX graph of itself
(\texttt{structural\_graph}), free edges are turned into labeled nodes so they
can participate in the matching, and two tensors are equal iff their graphs are
isomorphic (\texttt{is\_isomorphic}).
For speed, tensorgrad first compares Weisfeiler--Lehman hashes (and, when the
compiler backend is available, a refinement-based canonical fingerprint that is
computed compositionally and memoized), and only falls back to an exact
isomorphism search when the hashes collide.

Names make this subtler than plain graph isomorphism.
Inside a \texttt{Sum}, the terms $A$ and $A^T$ are isomorphic \emph{as graphs},
but $A + A^T$ must not simplify to $2A$: summands may only be merged under an
isomorphism that respects the sum's edge names.
Conversely, if the user declares
$A$ symmetric via \texttt{Variable.with\_symmetries}, the symmetry becomes part
of the graph, the renamed copy really is the same tensor, and $A + A^T = 2A$
follows automatically.
Isomorphism testing also powers \texttt{Product} simplification, where equal
subgraphs of a contraction are recognized and combined into powers
(Figure~\ref{fig:power}), and common-subexpression elimination during evaluation,
where a cached value for $A$ can be reused for $A^T$ by renaming its edges.

\begin{figure}[h]
\centering{
\def\svgwidth{\linewidth}
\import{figures/}{power.pdf_tex}
\caption{Combining equal parts of a product}
\label{fig:power}
}
\end{figure}

\section{Differentiation}

Derivatives follow Penrose's convention from the derivative chapters: taking
$\partial T/\partial X$ adds one new free edge to the diagram for every edge
of $X$.
Calling \texttt{t.grad(X)} wraps the expression in a \texttt{Derivative} node
and pushes it one step inward using the rules of this book, diagrams-as-data:
the derivative of a \texttt{Variable} with respect to itself becomes a product
of \texttt{Delta} tensors connecting old edges to new ones; the derivative of a
\texttt{Product} distributes by the product rule; a \texttt{Function} consults
its \texttt{FunctionSignature} for the chain rule from
Chapter~\ref{chapter:functions}.
Because a derivative is just another node in the graph, higher-order
derivatives like Hessians are simply \texttt{t.grad(x).grad(x)}, and
\texttt{Expectation} (over Gaussian variables, using the identities from the
statistics chapter) composes with everything else.

\section{Simplification and Evaluation}

Simplification is a set of local rewrite rules---evaluating deltas, flattening
nested sums and products, cancelling zeros, merging isomorphic terms---applied
in a depth-first pass by \texttt{simplify()}; \texttt{full\_simplify()} repeats
the pass until a fixed point is reached.
The result can be rendered back as a diagram with \texttt{to\_tikz} (that is how
many figures in this book were produced), evaluated numerically with
\texttt{evaluate}, which maps variables to PyTorch tensors and contracts
products with \texttt{einsum}, or compiled to a standalone PyTorch function
with \texttt{compile\_to\_callable} from \texttt{tensorgrad.extras.to\_pytorch}.

As a small end-to-end example, here is the gradient of the least-squares loss
from the machine-learning chapter:

\begin{lstlisting}
import tensorgrad as tg
import tensorgrad.functions as F

b, x, y = tg.symbols("b x y")
X = tg.Variable("X", b, x)
Y = tg.Variable("Y", b, y)
W = tg.Variable("W", x, y)

error = X @ W - Y            # contracts the shared edge x
loss = F.frobenius2(error)   # contract all edges with themselves
grad = loss.grad(W).full_simplify()
\end{lstlisting}

The resulting diagram is $2X^T(XW - Y)$, with no matrix-shaped reasoning
required from the user.
For installation instructions, an interactive playground, and many more
examples, see \url{https://github.com/thomasahle/tensorgrad}.
